% problem_id: S001-lemma-kunneth-single-sheaf
% source_id: S001 (Stacks Project)
% source_file: perfect.tex
% statement_locator: perfect.tex lines 5756--5761
% proof_locator: perfect.tex lines 5763--5975
% env: lemma
% id_note: label 在本来源内唯一
% pairing: tight (verified)
% status: complete (statement + author proof verbatim + reason + locators)
%
\begin{lemma}
\label{lemma-kunneth-single-sheaf}
In the situation above the map (\ref{equation-kunneth-single-sheaves}) is an
isomorphism if $S$ is affine, $\mathcal{F}$ and $\mathcal{G}$ are $S$-flat and
quasi-coherent and $X$ and $Y$ are quasi-compact with affine diagonal.
\end{lemma}

\begin{proof}
We strongly urge the reader to read the proof of
Varieties, Lemma \ref{varieties-lemma-kunneth} first.
Choose finite affine open coverings
$\mathcal{U} : X = \bigcup_{i \in I} U_i$ and
$\mathcal{V} : Y = \bigcup_{j \in J} V_j$.
This determines an affine open covering
$\mathcal{W} : X \times_S Y = \bigcup_{(i, j) \in I \times J} U_i \times_S V_j$.
Note that $\mathcal{W}$ is a refinement of
$\text{pr}_1^{-1}\mathcal{U}$ and of $\text{pr}_2^{-1}\mathcal{V}$.
Thus by the discussion in Cohomology, Section
\ref{cohomology-section-cech-cohomology-of-complexes}
we obtain maps
$$
\check{\mathcal{C}}^\bullet(\mathcal{U}, \mathcal{F})
\to
\check{\mathcal{C}}^\bullet(\mathcal{W}, p^*\mathcal{F})
\quad\text{and}\quad
\check{\mathcal{C}}^\bullet(\mathcal{V}, \mathcal{G})
\to
\check{\mathcal{C}}^\bullet(\mathcal{W}, q^*\mathcal{G})
$$
well defined up to homotopy and compatible with pullback maps on cohomology.
In Cohomology, Equation (\ref{cohomology-equation-needs-signs})
we have constructed a map of complexes
$$
\text{Tot}(
\check{\mathcal{C}}^\bullet(\mathcal{W}, p^*\mathcal{F})
\otimes_A
\check{\mathcal{C}}^\bullet(\mathcal{W}, q^*\mathcal{G}))
\longrightarrow
\check{\mathcal{C}}^\bullet(\mathcal{W},
p^*\mathcal{F} \otimes_{\mathcal{O}_{X \times_S Y}}
q^*\mathcal{G})
$$
which is compatible with the cup product on cohomology by
Cohomology, Lemma \ref{cohomology-lemma-diagrams-commute}.
Combining the above we obtain a map of complexes
\begin{equation}
\label{equation-kunneth-on-cech}
\text{Tot}(
\check{\mathcal{C}}^\bullet(\mathcal{U}, \mathcal{F})
\otimes_A
\check{\mathcal{C}}^\bullet(\mathcal{V}, \mathcal{G}))
\to
\check{\mathcal{C}}^\bullet(\mathcal{W},
p^*\mathcal{F}
\otimes_{\mathcal{O}_{X \times_S Y}}
q^*\mathcal{G})
\end{equation}
We claim this is the map in the statement of the lemma, i.e.,
the source and target of this arrow are the same as the source
and target of (\ref{equation-kunneth-single-sheaves}). Namely, by
Cohomology of Schemes, Lemma
\ref{coherent-lemma-quasi-coherent-affine-cohomology-zero}
and
Cohomology, Lemma \ref{cohomology-lemma-cech-complex-complex-computes}
the canonical maps
$$
\check{\mathcal{C}}^\bullet(\mathcal{U}, \mathcal{F})
\to
R\Gamma(X, \mathcal{F}),
\quad
\check{\mathcal{C}}^\bullet(\mathcal{V}, \mathcal{G})
\to
R\Gamma(Y, \mathcal{G})
$$
and
$$
\check{\mathcal{C}}^\bullet(\mathcal{W},
p^*\mathcal{F} \otimes_{\mathcal{O}_{X \times_S Y}} q^*\mathcal{G})
\to
R\Gamma(X \times_S Y,
p^*\mathcal{F} \otimes_{\mathcal{O}_{X \times_S Y}}
q^*\mathcal{G})
$$
are isomorphisms. On the other hand, the complex
$\check{\mathcal{C}}^\bullet(\mathcal{U}, \mathcal{F})$
is K-flat by Lemma \ref{lemma-K-flat} and we conclude that
$\text{Tot}(
\check{\mathcal{C}}^\bullet(\mathcal{U}, \mathcal{F})
\otimes_A
\check{\mathcal{C}}^\bullet(\mathcal{V}, \mathcal{G}))$
represents the derived tensor product
$R\Gamma(X, \mathcal{F}) \otimes_A^\mathbf{L} R\Gamma(Y, \mathcal{G})$
as claimed.

\medskip\noindent
We still have to show that (\ref{equation-kunneth-on-cech})
is a quasi-isomorphism. We will do this using dimension shifting.
Set $d(\mathcal{F}) = \max \{d \mid H^d(X, \mathcal{F}) \not = 0\}$.
Assume $d(\mathcal{F}) > 0$. Set $U = \coprod\nolimits_{i \in I} U_i$.
This is an affine scheme as $I$ is finite. Denote
$j : U \to X$ the morphism which is the inclusion $U_i \to X$
on each $U_i$. Since the diagonal of $X$ is affine, the morphism
$j$ is affine, see
Morphisms, Lemma \ref{morphisms-lemma-affine-permanence}.
It follows that $\mathcal{F}' = j_*j^*\mathcal{F}$ is $S$-flat, see
Morphisms, Lemma \ref{morphisms-lemma-pushforward-flat-affine}.
It also follows that $d(\mathcal{F}') = 0$ by combining
Cohomology of Schemes, Lemmas
\ref{coherent-lemma-relative-affine-cohomology} and
\ref{coherent-lemma-quasi-coherent-affine-cohomology-zero}.
For all $x \in X$ we have $\mathcal{F}_x  \to \mathcal{F}'_x$
is the inclusion of a direct summand: if $x \in U_i$,
then $\mathcal{F}' \to (U_i \to X)_*\mathcal{F}|_{U_i}$
gives a splitting. We conclude that
$\mathcal{F} \to \mathcal{F}'$ is injective and
$\mathcal{F}'' = \mathcal{F}'/\mathcal{F}$
is $S$-flat as well. The short exact sequence
$0 \to \mathcal{F} \to \mathcal{F}' \to \mathcal{F}'' \to 0$
of flat quasi-coherent $\mathcal{O}_X$-modules
produces a short exact sequence of complexes
$$
0 \to
\text{Tot}(
\check{\mathcal{C}}^\bullet(\mathcal{U}, \mathcal{F})
\otimes_A
\check{\mathcal{C}}^\bullet(\mathcal{V}, \mathcal{G})) \to
\text{Tot}(
\check{\mathcal{C}}^\bullet(\mathcal{U}, \mathcal{F}')
\otimes_A
\check{\mathcal{C}}^\bullet(\mathcal{V}, \mathcal{G})) \to
\text{Tot}(
\check{\mathcal{C}}^\bullet(\mathcal{U}, \mathcal{F}'')
\otimes_A
\check{\mathcal{C}}^\bullet(\mathcal{V}, \mathcal{G})) \to 0
$$
and a short exact sequence of complexes
$$
0 \to
\check{\mathcal{C}}^\bullet(\mathcal{W},
p^*\mathcal{F}
\otimes_{\mathcal{O}_{X \times_S Y}}
q^*\mathcal{G}) \to
\check{\mathcal{C}}^\bullet(\mathcal{W},
p^*\mathcal{F}'
\otimes_{\mathcal{O}_{X \times_S Y}}
q^*\mathcal{G}) \to
\check{\mathcal{C}}^\bullet(\mathcal{W},
p^*\mathcal{F}''
\otimes_{\mathcal{O}_{X \times_S Y}}
q^*\mathcal{G}) \to 0
$$
Moreover, the maps (\ref{equation-kunneth-on-cech}) between these are
compatible with these short exact sequences. Hence it suffices to prove
(\ref{equation-kunneth-on-cech})
is an isomorphism for $\mathcal{F}'$ and $\mathcal{F}''$. Finally,
we have $d(\mathcal{F}'') < d(\mathcal{F})$.
In this way we reduce to the case $d(\mathcal{F}) = 0$.

\medskip\noindent
Arguing in the same fashion for $\mathcal{G}$ we find that we
may assume that both $\mathcal{F}$ and $\mathcal{G}$
have nonzero cohomology only in degree $0$.
Observe that this means that $\Gamma(X, \mathcal{F})$
is quasi-isomorphic to the $K$-flat complex
$\check{\mathcal{C}}^\bullet(\mathcal{U}, \mathcal{F})$
of $A$-modules sitting in degrees $\geq 0$.
It follows that $\Gamma(X, \mathcal{F})$ is a flat $A$-module
(because we can compute higher Tor's against this module
by tensoring with the Cech complex).
Let $V \subset Y$ be an affine open. Consider the affine open covering
$\mathcal{U}_V : X \times_S V = \bigcup_{i \in I} U_i \times_S V$.
It is immediate that
$$
\check{\mathcal{C}}^\bullet(\mathcal{U}, \mathcal{F})
\otimes_A \mathcal{G}(V) =
\check{\mathcal{C}}^\bullet(\mathcal{U}_V,
p^*\mathcal{F} \otimes_{\mathcal{O}_{X \times Y}}
q^*\mathcal{G})
$$
(equality of complexes). By the flatness of $\mathcal{G}(V)$
over $A$ we see that
$\Gamma(X, \mathcal{F}) \otimes_A \mathcal{G}(V) \to
\check{\mathcal{C}}^\bullet(\mathcal{U}, \mathcal{F})
\otimes_A \mathcal{G}(V)$ is a quasi-isomorphism.
Since the sheafification of
$V \mapsto \check{\mathcal{C}}^\bullet(\mathcal{U}_V,
p^*\mathcal{F} \otimes_{\mathcal{O}_{X \times Y}}
q^*\mathcal{G})$ represents
$Rq_*(p^*\mathcal{F} \otimes_{\mathcal{O}_{X \times Y}} q^*\mathcal{G})$
by Cohomology of Schemes, Lemma
\ref{coherent-lemma-separated-case-relative-cech}
we conclude that
$$
Rq_*(p^*\mathcal{F} \otimes_{\mathcal{O}_{X \times Y}} q^*\mathcal{G})
\cong
\Gamma(X, \mathcal{F}) \otimes_A \mathcal{G}
$$
on $Y$ where the notation on the right hand side indicates the module
$$
b^*\widetilde{\Gamma(X, \mathcal{F})} \otimes_{\mathcal{O}_Y} \mathcal{G}
$$
Using the Leray spectral sequence for $q$ we find
$$
H^n(X \times_S Y, p^*\mathcal{F} \otimes_{\mathcal{O}_{X \times Y}}
q^*\mathcal{G}) =
H^n(Y,
b^*\widetilde{\Gamma(X, \mathcal{F})} \otimes_{\mathcal{O}_Y} \mathcal{G})
$$
Using Lemma \ref{lemma-cohomology-base-change} for the morphism
$b : Y \to S = \Spec(A)$ and using that $\Gamma(X, \mathcal{F})$
is $A$-flat we conclude that
$H^n(X \times_S Y, p^*\mathcal{F} \otimes_{\mathcal{O}_{X \times Y}}
q^*\mathcal{G})$ is zero for $n > 0$ and isomorphic to
$H^0(X, \mathcal{F}) \otimes_A H^0(Y, \mathcal{G})$ for $n = 0$.
Of course, here we also use that $\mathcal{G}$ only has
cohomology in degree $0$.
This finishes the proof (except that we should check that the
isomorphism is indeed given by cup product in degree $0$; we omit
the verification).
\end{proof}
